% ===== 原文 pp.40--52 =====
% 译者 tr-03；§5 第四讲（pp.40--45）、§6 第五讲（pp.46--52）
% 页图 p040.png--p052.png 为公式唯一真值；文字层仅用于定位。

% [原文 p.40]
\section{第四讲（Chern 特征与 Chern 类）}

我们已经看到，$K^n(X,A)$ 是一个可乘的广义上同调理论。此外我们还辨认出了映射
\[ K^n(X,A)\times K^m(Y,B)\to K^{n+m}(X\times Y,\, X\times B\cup A\times Y). \]

成立的是
\[ K^*(X,A)\otimes\Q \]
是一个普通上同调理论。

我们假定 $X$ 是「好」（nice）的拓扑空间（例如有限单纯复形或有限 CW 复形）。

于是我们定义：
\[ H^{ev}(X;\Q) = \bigoplus_{n\ge 0} H^{2n}(X;\Q) \]
以及
\[ H^{odd}(X;\Q) = \bigoplus_{n\ge 0} H^{2n+1}(X;\Q). \]

于是有如下基本结果：

\begin{thm}{1}
存在自然同构
\[ \ch : K^0(X)\otimes\Q\to H^{ev}(X;\Q) \]
以及
\[ \ch : K^1(X)\otimes\Q\to H^{dd}(X;\Q) \]
它们保持乘积。这些同构称为 \emph{Chern 特征}。
\end{thm}
% [Tutor Note: p.40 定理 1 第二个展示式右端原文印作 $H^{dd}$，应为 $H^{odd}$：本讲开头定义的正是 $H^{odd}(X;\Q)=\bigoplus_{n\ge0}H^{2n+1}(X;\Q)$，且 p.46 复述同一结论时印的是 $H^{odd}(X;\Q)$。此处照原样译作 $H^{dd}$，不改动。]

\begin{proof}
本讲末尾给出了该证明的概要。
\end{proof}
% [Tutor Note: p.40 原文 "There is a scetch of the proof" 中 scetch 系 sketch 之误；译文按原意作「概要」。另：原文把 $K^0(X)$、$K^1(X)$ 压排作 $K^0X$、$K^1X$，译文按项目体例统一写作带括号的形式，符号本身未变。]

继续往下，我们需要复线丛的示性类。设 $L\to X$ 是复线丛（下文我们将直接把它记作 $L_X$），并令
\[ \Vect_1 X = [X, G_1], \]
% [原文 p.41]
其中
\[ f_L : X\to G_1 = \bigcup_n G_1(\C^n) = \bigcup_n \C P^{n-1} := \C P^\infty. \]
于是
\[ H^*(\C P^\infty;\Z) = \Z[u], \]
其中
\[ u\in H^2(\C P^\infty;\Z). \]

\begin{defn}{1}
$L$ 的第一 \emph{Chern} 类为
\[ c_1 L = f_L^*(u)\in H^2(X;\Z). \]
\end{defn}

\begin{lem}{1}
映射
\[ c_1 : \Vect_1(X)\to H^2(X;\Z) \]
是群同构。
\end{lem}
% [Tutor Note: p.41 引理 1 中的 $\Vect_1(X)$ 必须理解为「复线丛的同构类在张量积下构成的群」——其逆元由 $L^*$ 给出（见下两段的 $L^*\otimes L\simeq\C_X$），单位元是 $\C_X$。若只把 $\Vect_1(X)$ 看成同构类的集合，「群同构」便无从谈起。核对点：$c_1(L^*)=-c_1(L)$ 正是这条群结构的推论。]

给定两个线丛 $L_1$ 与 $L_2$，可以作它们的张量积 $L_1\otimes L_2$，其中 $\C_X$ 是张量积 $\otimes$ 的单位元。记
\[ L_X^* := \Hom(L_X,\C) \]
为对偶复线丛。则
\[ L^*\otimes L\simeq \C_X. \]
此外
\[ c_1(L_1\otimes L_2) = c_1 L_1 + c_2 L_2. \]
% [Tutor Note: p.41 上式右端原文印作 $c_2L_2$，应为 $c_1L_2$（已用 500 dpi 裁图复核，原文确印 $c_2$）。判断依据：$L_2$ 是线丛，次数大于 1 的 Chern 类全为零，若照原式则等式一般不成立；正确形式是 $c_1(L_1\otimes L_2)=c_1L_1+c_1L_2$，来自 $\det(L_1\otimes L_2)\cong\det L_1\otimes\det L_2$。正文照原样译。]
% [Tutor Note: p.41 同段原文 "the duall complex line bundle" 中 duall 系 dual 之误（800 dpi 裁图复核：确为 duall）；译文作「对偶」，不保留多出的那个 l。]

给定 $n$ 维向量丛 $\xi\to X$，我们考虑丛 $\pi : \mathbb{P}(\xi)\to X$，其纤维为 $\pi^{-1}(x) = \mathbb{P}(F_x(\xi))\simeq \C P^{n-1}$。诱导映射
\[ \pi^* : H^*(X)\to H^*(\mathbb{P}(\xi)) \]
是环同态，它使 $H^*(\mathbb{P}(\xi))$ 成为环 $H^*(X)$ 上的模（整系数）。

接下来考虑映射
\[ H^*(X)\otimes H^*(\mathbb{P}(\xi))\to H^*(\mathbb{P}(\xi)), \]
% [原文 p.42]
其中
\[ x\otimes y\mapsto \pi^*(x)y. \]

在 $\mathbb{P}(\xi)$ 上存在线丛 $L_\xi$，使得 $1\in H^0(\mathbb{P}(\xi))$、$c_1(L_\xi) = u\in H^2(\mathbb{P}(\xi))$；我们再作 $u^2$、$u^3$ 等等。于是有下述结论：

\begin{thm}{2}
$H^*(\mathbb{P}(\xi))$ 是 $H^*(X)$ 上的自由模，基为 $1,u,u^2,\ldots u^{n-1}$。
\end{thm}

\begin{proof}
我们描述基本想法。

情形 1：$\xi = X\times\C^1$，此时 $\mathbb{P}(\xi) = X\times\C P^{n-1}$，而
\[ H^*(X\times\C P^{n-1}) = H^*(X)\otimes H^*(\C P^{n-1}) \]
可用代数拓扑中的 Kunneth 公式算出。为完成证明，我们使用 Mayer-Vietories 归纳论证。
\end{proof}
% [Tutor Note: p.42 上句两处拼写：Kunneth 应为 Künneth（Künneth 公式），Mayer-Vietories 应为 Mayer--Vietoris（Mayer--Vietoris 序列）。正文照原样保留，便于对照原文复核。]

\begin{defn}{2}
（\emph{Chern 类}）设 $\xi\to X$ 是向量丛，作 $\mathbb{P}(\xi)$ 并取类 $u\in H^2(\mathbb{P}(\xi);\Z)$。则
\[ -u^n = c_1(\xi)u^{n-1}+c_2(\xi)u^{n-2}+\ldots+c_k(\xi)u^{n-k}+\ldots+c_n(\xi). \]
上式定义了 Chern 类 $c_i(\xi)\in H^{2i}(X;\Z)$。
\end{defn}
% [Tutor Note: p.42 定义 2 的等式左端带负号，这是 $u=c_1(L_\xi)$ 这一符号约定所致：移项即得 $u^n+c_1(\xi)u^{n-1}+\ldots+c_n(\xi)=0$，即 $u$ 是 $c(\xi)(u)$ 的根。核对点：若改用 $u=-c_1(L_\xi)$，同一等式会写成 $u^n=c_1u^{n-1}-\ldots$，各项符号全变，故抄录时不可略去这个负号。]

\para{问题：}验证 Chern 类的下列性质：
\begin{enumerate}[label=\arabic*.]
\item （拉回）$c_i(f^*(\xi)) = f^*(c_i(\xi))$。
\item （\emph{Whitney 和公式}）
\[ c(\xi) = 1+c_1(\xi)+c_2(\xi)+\ldots+c_n(\xi), \]
\[ c(\xi\oplus\eta) = c(\xi)c(\eta) \]
以及
\[ c_k(\xi\oplus\eta) = \sum_{i+j=k}c_i(\xi)c_j(\eta). \]
\end{enumerate}
% [Tutor Note: p.42 「问题」第 2 条里同时出现了总 Chern 类 $c(\xi)$ 与分量 $c_k(\xi)$：前者是形式和 $1+c_1+\ldots+c_n$，故 $c(\xi\oplus\eta)=c(\xi)c(\eta)$ 是形式幂级数的乘法；后者按次数拆开即得 $c_k(\xi\oplus\eta)=\sum_{i+j=k}c_i(\xi)c_j(\eta)$。二者是同一条 Whitney 公式的两种写法，核对方法：把前一式按 $k$ 次分量比较系数。]

\begin{lem}{2}
（分裂原理）设 $\xi\to X$ 是复向量丛。则存在空间 $F(\xi)$ 以及映射 $f : F(\xi)\to X$，使得：
\begin{enumerate}[label=\roman*.]
\item $f^* : H^*(X)\to H^*(F(\xi))$ 是单射；并且
\item $f^*(\xi)\simeq L_1\oplus L_2\oplus\ldots\oplus L_n$，其中 $n$ 是 $\xi$ 的维数。
\end{enumerate}
（注意同一结论对 K 理论也成立，即 $f^* : K^*(X)\to K^*(F(\xi))$ 也是单射）。
\end{lem}

\begin{proof}
由上面的引理 1，$f_1 : \mathbb{P}(\xi)\to X$ 是单射。此外
\[ L_\xi\subset f^*(\xi)\Rightarrow f^*(\xi) = L_\xi\oplus \tilde\xi, \]
而要完成证明，我们使用归纳法。
% [原文 p.43]
\end{proof}
% [Tutor Note: p.42 引理 2 的证明开头「$f_1:\mathbb{P}(\xi)\to X$ 是单射」有误：$\mathbb{P}(\xi)$ 的纤维是 $\C P^{n-1}$，$f_1$ 不可能单射；原意应是 $f_1^*:H^*(X)\to H^*(\mathbb{P}(\xi))$ 单射，而这正是定理 2 的内容（不是引理 1）。下文 $f^*(\xi)=L_\xi\oplus\tilde\xi$ 用的也仍是「$f_1^*$ 单射」这一性质。正文照原样译，读者按 $f_1^*$ 理解即可。]
% [Tutor Note: p.43 接上一句，原文跨页断词作 "induc-/tion"（即 induction），译文把它与 p.42 末句合并译为「使用归纳法」；方框 □ 落在 p.43 顶端，与原文一致。]

接下来我们想构造 Chern 特征 $\ch : K^0(X)\to H^{ev}(X;\Q)$，使它满足所要求的性质（见上面的「问题」；为简化记号，我们略去方括号 $[\xi]$）：
\begin{itemize}
\item $\ch(\xi\oplus\eta) = \ch(\xi)+\ch(\eta)$
\item $\ch(\xi\otimes\eta) = \ch(\xi)\ch(\eta)$.
\end{itemize}
% [Tutor Note: p.43 原文此处写作 "Chen character"，Chen 系 Chern 之误（同段后文与 p.40 均作 Chern）；译名按「Chern 特征」处理。同段 "we ommit" 中 ommit 亦系 omit 之误。]

线丛 $L$ 的 $\ch(L)$ 该是什么？

设 $c_1(L) = x$、$c_1(\tilde L) = y$。则 $c_1(L\otimes\tilde L) = x+y$。这使人联想到幂级数 $F(x)$（把 $\ch(L) = F(x)$ 设想成它），其中 $F(x+y) = F(x)F(y)$，从而
\[ F(x) = e^x. \]
于是我们得到下述定义：

\begin{defn}{3}
对复线丛 $L$，定义
\[ \ch(L) = e^{c_1(L)} = 1+c_1(L)+\frac{c_1^2(L)}{2}+\ldots \]
\end{defn}
% [Tutor Note: p.43 定义 3 的依据是：在形式幂级数环里，满足 $F(x+y)=F(x)F(y)$ 且 $F(0)=1$ 的 $F$ 唯一，就是 $e^x$。这一步同时说明了为什么 $\ch$ 必须取有理系数（出现分母 $k!$），从而定理 1 里要作 $K^0(X)\otimes\Q$——整数系数的 K 理论本身给不出 $\ch$。]

对任意的复向量丛 $E$，我们有
\[ E = L_1\oplus L_2\oplus\ldots\oplus L_n \]
其中 $c_1(L_1) = x_1$、$c_1(L_2) = x_2$、……、$c_1(L_n) = x_n$。则
\[ \ch(E) = e^{x_1}+e^{x_2}+\ldots+e^{x_n} \]
从而
\[ \resizebox{0.98\textwidth}{!}{$\displaystyle \ch(E) = c(L_1)c(L_2)\ldots c(L_n) = (1+x_1)(1+x_2)\ldots(1+x_n) = 1+c_1(E)+c_2(E)+\ldots+c_n(E)$} \]
这是因为 $c(L_i) = 1+x_i$。

于是
\[ c_i(E) = \sigma_i(x_1,x_2,\ldots x_n) \]
其中 $\sigma_i$ 是第 $i$ 个初等对称函数。

% [原文 p.44]
例如 $c_1(E) = x_1+x_2+\ldots+x_n$，而 $c_2(E) = \prod_{i<j}x_ix_j$，从而
\[ e^{x_1}+e^{x_2}+\ldots+e^{x_n} = n+S_1(\sigma_1,\sigma_2,\ldots,\sigma_n)+S_2(\sigma_1,\sigma_2,\ldots,\sigma_n)+\ldots \]

\begin{defn}{4}
\[ \ch(\xi) = \dim\xi + S_1(c_1(\xi),c_2(\xi),\ldots,c_n(\xi))+S_2(c_1(\xi),c_2(\xi),\ldots,c_n(\xi))+\ldots \]
\end{defn}
% [Tutor Note: p.44 定义 4 中的 $S_i$ 原文未加定义，它就是把幂和用初等对称函数表出的第 $i$ 个 Newton 和多项式：$S_1=\sigma_1$、$S_2=\sigma_1^2-2\sigma_2$（正是本页紧接着算的例子）。到 p.50 同一对象改记作 $\Phi^k$（幂和）并与 $\Psi^k$ 对应，故读 p.50 时应把 $S_i$ 与 $\Phi^i$ 认作同一物。]

作为例子，我们来算出 $S_2$ 的情形：
\[ x_1^2+x_2^2+\ldots+x_n^2 = (x_1+x_2+\ldots+x_n)^2-2\prod_{i<j}(x_ix_j) = \sigma_1^2-2\sigma_2. \]
% [Tutor Note: p.44 本页两处「$\prod_{i<j}$」原文都印成大写 $\Pi$（500 dpi 裁图复核，确系 $\Pi$）：$c_2(E)=\prod_{i<j}x_ix_j$ 与 $2\prod_{i<j}(x_ix_j)$；按定义，$c_2$ 是初等对称函数 $\sigma_2=\sum_{i<j}x_ix_j$，第二式也应为 $(\sum x_i)^2-2\sum_{i<j}x_ix_j$。若照字面取乘积，$n\ge3$ 时两式均不成立（$n=2$ 时只有一对 $i<j$，乘积与和恰好相同，故不易察觉）。正文照原文保留 $\prod$。]
于是
\[ \ch(\xi) = \dim\xi+c_1(\xi)+\frac{1}{2}[c_1^2(\xi)-2c_2(\xi)]+\ldots \]

回忆所期望的两条性质 $\ch(\xi\oplus\eta) = \ch(\xi)+\ch(\eta)$ 与 $\ch(\xi\otimes\eta) = \ch(\xi)\ch(\eta)$ 都来自分裂原理。

我们定义了
\[ \ch : K^0(X)\to H^{ev}(X;\Q). \]

\begin{defn}{5}
\[ \ch([\xi]-[\eta]) = \ch(\xi)-\ch(\eta). \]
\end{defn}

于是
\begin{center}
\begin{tikzcd}
K^1(X) \arrow[r] \arrow[d,equal] & H^{odd}(X;\Q) \arrow[d,equal] \\
\wt K^0(SX) \arrow[r,"\ch"'] & \wt H^{ev}(SX;\Q)
\end{tikzcd}
\end{center}
% [Tutor Note: p.44 图中两条竖线都是「等号」而非同构箭头，表示上下两行被同一对象等同：$K^1(X)\cong\wt K^0(SX)$（悬挂/约化 K 的 Bott 描述）与 $H^{odd}(X;\Q)\cong\wt H^{ev}(SX;\Q)$（上同调悬挂）。因此图中「$\ch$ 与悬挂同构相容」这一断言才说得通；抄图时若把等号改成箭头，含义就变了。]

我们终于来到本讲开头所许诺的定理 1 的证明：

\begin{proof}[定理 1 的证明]
% [原文 p.45]
第 1 步。若 $X = S^2$，则 $\wt K^0(S^2) = \Z$ 且 $K^0(S^2) = \Z\oplus\Z$。此外 $\ch(1) = 1$，而 $\ch(\xi) = 1+u$，因为 $H^{ev}(S^2;\Q) = \Q\oplus\Q$（$1$ 对应第一个因子 $\Q$，$u$ 对应第二个）。

第 2 步。当 $X = S^{2n}$ 时有同构，因为 Chern 特征与乘积可交换。

第 3 步。当 $X = S^{2n+1}$ 时有同构（直接计算）。

第 4 步。设 $X$ 是有 $N$ 个单形的单纯复形，$A\subset X$ 由去掉一个最高维单形得到。但此时 $X/A = S^n$，我们使用正合列与五引理。
\end{proof}
% [Tutor Note: p.45 第 1 步的 $\ch(\xi)=1+u$ 可自行核对：$\xi$ 取 $S^2$ 上生成 $K^0(S^2)$ 的线丛（即 Hopf 丛 $J$），$c_1(J)=u$，由定义 3 得 $\ch(J)=e^{u}=1+u+\frac{u^2}{2}+\ldots$；而 $u\in H^2(S^2)$ 故 $u^2=0$，级数截断为 $1+u$。另一核对点：约化侧 $\ch([J]-1)=u$，正与 $\wt K^0(S^2)=\Z$ 的生成元对应。]

% [原文 p.46]
\section{第五讲（K 理论中的运算、对称积与 Adams 运算）}

上一节我们定义了 Chern 特征：
\[ \ch : K^0(X)\to H^{ev}(X;\Q) \]
以及
\[ \ch : K^1(X)\to H^{odd}(X;\Q) \]
并且叙述了基本结果：映射
\[ \ch : K^*(X)\otimes\Q\to H^*(X;\Q) \]
是同构。

现在我们来研究 \emph{K 理论中的运算}。

设 $V$ 是有限维复向量空间，维数设为 $n$。由线性代数我们知道外幂
\[ \wedge^k V = (V\otimes V\otimes\ldots\otimes V)/(x_1\otimes x_2\otimes\ldots\otimes x_k - sign(\sigma)\,x\otimes_{\sigma(1)}x\otimes_{\sigma(2)}\ldots\otimes_{\sigma(k)}x). \]
其中张量积取了 $k$ 个因子。显然 $\wedge^0V = \C$，$\wedge^1V = V$，……，$\wedge^kV$，……，$\wedge^nV$，其中 $\dim\wedge^nV = 1$，更一般地 $\dim\wedge^kV = (n!)/[k!(n-k)!]$。
% [Tutor Note: p.46 上式的置换项原文排作 $x\otimes_{\sigma(1)}x\otimes_{\sigma(2)}\ldots\otimes_{\sigma(k)}x$（把下标挂到了 $\otimes$ 上），正确形式应为 $x_{\sigma(1)}\otimes x_{\sigma(2)}\otimes\ldots\otimes x_{\sigma(k)}$。判断依据：$\wedge^kV$ 是 $V^{\otimes k}$ 对交错化作用的商，被等同掉的正是 $x_1\otimes\ldots\otimes x_k-\sgn(\sigma)x_{\sigma(1)}\otimes\ldots\otimes x_{\sigma(k)}$；取 $k=2$、$\sigma=(12)$ 代入正确形式即得 $x_1\otimes x_2-x_2\otimes x_1$，正是熟知的反对称关系。正文照原样译。]

我们可以把上述构造以直接的方式推广到某空间 $X$ 上的向量丛 $\xi$：作外幂 $\wedge^k\xi$，它也是 $X$ 上的向量丛，纤维为
\[ F_x(\wedge^k\xi) = \wedge^k(F_x(\xi)). \]

我们想构造映射
\[ \wedge^k : \Vect(\xi)\to K^0(X). \]
% [Tutor Note: p.46 上式中的 $\Vect(\xi)$ 按上下文应为 $\Vect(X)$（$X$ 上向量丛同构类的集合）：$\wedge^k$ 是对丛取外幂，自变量应是丛而不是 $\xi$ 这个丛变量。依据是 p.47 同一映射写作 $\wedge_t:\Vect(X)\to K^0(X)[[t]]$。正文照原样译。]

若 $V,W$ 是两个向量空间，则
\[ \wedge^k(V\oplus W) = \sum_{i+j=k}\wedge^i V\otimes\wedge^j W \]
% [原文 p.47]
相应的向量丛版本为
\[ \wedge^k(\xi\oplus\eta) = \sum_{i+j=k}\wedge^i(\xi)\otimes\wedge^j(\eta). \]
注意与 Chern 类的类比：
\[ c_k : \Vect(X)\to H^{ev}(X;\Q) \]
其中
\[ c_k(\xi\oplus\eta) = \sum_{i+j=k}c_i(\xi)c_j(\eta). \]
% [Tutor Note: p.47 上式中 $c_k$ 的靶按 p.42 定义 2 应为 $H^{2k}(X;\Z)$；原文写作 $H^{ev}$，可视为把 $c_k$ 直接看成偶次部分的元素（$H^{2k}\subset H^{ev}$）。真正的类比点是：$c_k$ 与 $\wedge^k$ 都满足同一条「加法变卷积」的乘性公式。]

给定向量丛 $\xi$，我们可以构造 $\wedge_t(\xi)\in K^0(X)[[t]]$，它是 $t$ 的形式幂级数、系数取自群 $K^0(X)$，换言之
\[ \wedge_t\xi = \sum_{k=0}^\infty t^k\wedge^k\xi. \]
于是
\[ \wedge^k(\xi\oplus\eta) = \sum_{i+j=k}\wedge^i(\xi)\otimes\wedge^j(\eta) \]
变成
\[ \wedge_t(\xi\oplus\eta) = \wedge_t(\xi)\wedge_t(\eta) \]
这意味着我们有映射
\[ \wedge_t : \Vect(X)\to K^0(X)[[t]]. \]
我们的约定是 $\wedge^0(\xi) = 1\in K^0(X)$，因此 $\wedge_t(\xi)$ 是首项为 1 的形式幂级数。从而 $\wedge_t(\xi)$ 是 $K^0(X)[[t]]$ 中的可逆元。

记 $G(K^0(X)[[t]])$ 为环 $K^0(X)[[t]]$ 的单位群，即可逆元构成的群。于是得到如下交换图
\begin{center}
\begin{tikzcd}
\Vect(X) \arrow[r,"\wedge_t"] \arrow[d] & G(K^0(X)[[t]]) \\
K^0(X) \arrow[r,equal] & K^0(X) \arrow[u]
\end{tikzcd}
\end{center}
% [Tutor Note: p.47 图中底边的双线仍是等号：左右两个 $K^0(X)$ 指同一个群，即下路的终点与上路的起点是同一对象。核对方法：$\wedge_t(\xi)$ 的常数项恒为 $1$，故 $\wedge_t$ 的像落在单位群 $G(K^0(X)[[t]])$ 里而不是整个形式幂级数环里——这正是这张图要表达的意思；若把 $\wedge_t$ 看成映入 $K^0(X)[[t]]$ 的映射，图的右上方就该改写成整个环。]

% [原文 p.48]
给定 $x\in K^0(X)$，我们定义 $\wedge^i(x)$ 为 $\wedge^i_t(x)$ 中 $t^i$ 的系数。若 $x = [\xi]-[\eta]$，则 $\wedge_t(x) = \wedge_t(\xi)\wedge_t(\eta)^{-1}$。
% [Tutor Note: p.48 原文把 $\wedge^i_t(x)$ 的上下标都排了出来（上标 $i$、下标 $t$），指的是 $\wedge_t(x)$ 的 $t^i$ 系数；同一个符号 $\wedge^i(x)$ 在 p.49 又说成映射 $\wedge^i:K^0(X)\to K^0(X)$，两处是同一件事（先算出 $\wedge_t$ 再取系数），并不矛盾。抄录时注意别把上标 $i$ 丢掉。]

例如，若 $L$ 是线丛，则 $\wedge_t(L) = 1+tL$，且
\[ \wedge_{-t}(-L) = \frac{1}{1-tL} = 1+tL+t^2L^2+\ldots \]

这里我们要提一句：为 $t=$ something 取定一个值难以说通（这正是我们在谈形式幂级数的原因），但系数却完全有意义（但愿如此）。
% [Tutor Note: p.48 原文作 "for $t=$something"，等号后紧接字母、无空格（页图复核）；译文按「$t=$ something」排。这句的用意是：$\wedge_t$ 只是形式幂级数，$t$ 不取数值，但 $t^k$ 的系数（即 $\wedge^k$）有确定含义。]

接下来回忆线性代数中的对称幂
\[ S^k(V) = (V\otimes V\otimes\ldots\otimes V)^{\Sigma_k} \]
% [Tutor Note: p.48 上式右上标 $\Sigma_k$ 表示取 $\Sigma_k$-不变量（对称群作用下的不动子空间），即对称张量构成的子空间；这是「$S^k$ 是 $\wedge^k$ 的对称类比」的确切含义。核对点：$\dim S^kV=\binom{n+k-1}{k}$，而 $\dim\wedge^kV=\binom{n}{k}$，两者在 $k=2$ 时分别为 $\frac{n(n+1)}{2}$ 与 $\frac{n(n-1)}{2}$，之和为 $n^2=\dim(V\otimes V)$。]
其中张量积假定有 $k$ 个因子。类似地可以作向量丛的对称幂 $S^k(\xi)$，然后定义
\[ S_t(\xi) = \sum_{k=0}^\infty t^k S^k(\xi) \]
我们的约定是 $S^0(\xi) = 1$。此外，由线性代数有
\[ S^k(V\oplus W) = \sum_{i+j=k}S^iV\otimes S^jW. \]
用同样的方法可以得到如下交换图
\begin{center}
\begin{tikzcd}
\Vect(X) \arrow[r,"S_t"] \arrow[d] & G(K^0(X)[[t]]) \\
K^0(X) \arrow[r,equal] & K^0(X) \arrow[u]
\end{tikzcd}
\end{center}

% [原文 p.49]
例如，若 $L$ 是线丛，则
\[ S_t(L) = 1+tL+t^2L^2+\ldots \]
这是因为
\[ S^kL = L^{\otimes k}. \]
因此
\[ S_t(L) = \wedge_{-t}(-L), \]
从而
\[ \wedge_{-t}(L)S_t(L) = 1 \Leftrightarrow \wedge_t(L)S_{-t}(L) = 1\Rightarrow \wedge_t(\xi)S_{-t}(\xi) = 1 \]
其中任意的向量丛 $\xi$ 可以写成线丛之和。

由 K 理论的分裂原理推得，对任意向量丛 $\xi$ 都有
\[ \wedge_t(\xi)S_{-t}(\xi) = 1 \]
于是
\[ \wedge_t([\xi]-[\eta]) = (\wedge_t\xi)(S_{-t}\eta) \]
并且
\[ \wedge^k([\xi]-[\eta]) = \sum_{i+j=k}(-1)^j\wedge^i(\xi)S^j(\eta). \]
% [Tutor Note: p.49 最后这条式子的符号来自 $\wedge_t(\xi)S_{-t}(\xi)=1$：由 $\wedge_t([\xi]-[\eta])=\wedge_t\xi\cdot\wedge_t(\eta)^{-1}$ 与 $\wedge_t(\eta)^{-1}=S_{-t}(\eta)$ 得 $(\wedge_t\xi)(S_{-t}\eta)$，再按 $t^k$ 取系数时 $S_{-t}\eta$ 的每个 $t^j$ 项带 $(-1)^j$。核对点：取 $\eta=0$ 应退化为 $\wedge^k(\xi)$，此时求和只剩 $j=0$ 项，与原文一致。]

我们作一个重要注记：若 $f : X\to Y$ 是映射，则对某个 $x\in K^0(Y)$ 有
\[ f^*(\wedge^i(x)) = \wedge^i(f^*x) \]
其中 $\wedge^i : K^0(X)\to K^0(X)$。
% [Tutor Note: p.49 这条注记须按字面读：$x$ 取自 $K^0(Y)$，$f^*x\in K^0(X)$；而 $\wedge^i$ 的靶写的是 $K^0(X)\to K^0(X)$，所以 $\wedge^i(f^*x)$ 是先拉回再作用。原文把 $\wedge^i$ 的定义域写在句子末尾，容易误读成 $x\in K^0(X)$；一旦读错，「$f^*$ 与 $\wedge^i$ 可交换」就无从对照。]

现在我们来研究 K 理论中的 \emph{Adams 运算}。

设 $x\in K_0(X)$ 且 $\Psi_t(x)\in K^0(X)[[t]]$。
% [Tutor Note: p.49 原文此处把 $K^0(X)$ 排成了下标形式 $K_0(X)$（同句后半即作 $K^0(X)$，本讲其余各处亦均用上标）；按上下文应读作 $K^0(X)$。正文照原文保留。]

我们定义
\[ \Psi_t(x) := \dim(x)-\frac{t}{\wedge_{-t}(x)}\frac{d}{dt}\wedge_{-t}(x) = \dim(x)-t\frac{d}{dt}\log(\wedge_t(x)). \]
% [Tutor Note: p.49 末式右端第二个等号处，$\log$ 里原文印作 $\wedge_t(x)$，按左端的定义应为 $\wedge_{-t}(x)$。判断依据：把线丛 $L$ 代入左端（$\wedge_{-t}(L)=1-tL$）得 $\Psi_t(L)=1+\frac{tL}{1-tL}$，这正是 p.51 证明第 2 条算出的值；若改用 $\log(\wedge_t(L))=\log(1+tL)$，得到的是 $1-\frac{tL}{1+tL}$，与 p.51 不符。]

% [原文 p.50]
所有系数都是 $K^0(X)$ 中元素的整倍数。
% [Tutor Note: p.50 原文 "integral multiplets" 中 multiplets 系 multiples 之误；译文作「整倍数」。]

此外我们还通过方程
\[ \sum t^k\Psi^k(x) = \Psi_t(x) \]
定义量 $\Psi^k(x)$。

\begin{defn}{1}
映射
\[ \Psi^k : K^0(X)\to K^0(X) \]
称为 K 理论中的 \emph{Adams 运算}。
\end{defn}

\para{例：}
\[ \Psi^1-\wedge^1 = 0 \]
\[ \Psi^2-\Psi^1\wedge^1+2\wedge^2 = 0 \]
\[ \Psi^3-\Psi^2\wedge^1+\Psi^1\wedge^2-3\wedge^2 = 0 \]
一般地
\[ \Psi^k-\Psi^{k-1}\wedge^1\pm\ldots\pm k\wedge^2 = 0. \]
% [Tutor Note: p.50 例中第三行末项原文印作 $3\wedge^2$，应为 $3\wedge^3$；紧接着的一般式末项同样印作 $k\wedge^2$，应为 $k\wedge^k$（两处均已 500 dpi 裁图复核，确系原印）。判断依据：这三行是 Newton 恒等式 $\Psi^k-\Psi^{k-1}\wedge^1+\Psi^{k-2}\wedge^2-\ldots+(-1)^{k-1}k\wedge^k=0$ 在 $k=1,2,3$ 的特例——符号交替、末项系数为 $k$、幂次为 $k$。旁证：本页下一段说 $\Psi^k\leftrightarrow\Phi^k$、$\wedge^i\leftrightarrow\sigma_i$（幂和与初等对称函数），而 Newton 恒等式中 $\sigma_k$ 前的系数恰为 $(-1)^{k-1}k$。]

Adams 运算联系幂和 $\Phi^k(x_1,x_2,\ldots,x_n) = x_1^k+x_2^k+\ldots+x_n^k$ 与初等对称函数 $\sigma_i(x_1,x_2,\ldots,x_n)$，即 $\Psi^k\leftrightarrow\Phi^k$ 以及 $\wedge^i\leftrightarrow\sigma_i$。

我们把 Adams 运算的基本性质总结如下：

\begin{prop}{1}
Adams 运算具有下列性质：
\begin{enumerate}[label=\arabic*.]
\item \[ \Psi^k(x+y) = \Psi^k(x)+\Psi^k(y). \]
\item 对线丛 $L$，
\[ \Psi^k(L) = L^k \]
\item \[ \Psi^k(xy) = \Psi^k(x)\Psi^k(y) \]
% [原文 p.51]
\item \[ \Psi^k(\Psi^l(x)) = \Psi^{kl}(x). \]
\item 若 $p$ 是素数，则
\[ \Psi^p(x) = x^p+py = x^p modp \]
\item 若 $u\in\wt K^0(S^{2n})$，则
\[ \Psi^k(u) = k^n u. \]
\end{enumerate}
\end{prop}
% [Tutor Note: p.51 第 5 条里的 $y$ 原文未作任何说明，应理解为某个元素 $y\in K^0(X)$，于是整条结论即 $\Psi^p(x)\equiv x^p\pmod p$（差落在 $pK^0(X)$ 中）。这也是 Frobenius 同余的 K 理论版本：$p$ 为素数时 $(L_1+\ldots+L_n)^p\equiv L_1^p+\ldots+L_n^p$。核对点：p.51 末段正是取 $x=L_1+\ldots+L_n$ 把它重算了一遍。另：原文 "$x^pmodp$" 中 modp 连排、无空格，译文按原样作 $x^p modp$。]

\begin{proof}
我们来证明其中几条。对第一条，我们有
\[ \resizebox{0.98\textwidth}{!}{$\displaystyle \Psi_t(x+y) = \dim(x+y)-\frac{t}{\wedge_{-t}(x+y)}\frac{d}{dt}\wedge_{-t}(x+y) = \dim(x+y)-\frac{t}{\wedge_{-t}(x)\wedge_{-t}(y)}\times$} \]
\[ \times \frac{d}{dt}(\wedge_t(x)\wedge_{-t}(y)) \]
然后使用 Leibniz 法则即得所要的结果 $\Psi_t(x)+\Psi_t(y)$。
% [Tutor Note: p.51 第二个等号把 $\wedge_{-t}(x+y)$ 拆成 $\wedge_{-t}(x)\wedge_{-t}(y)$（这一步用的是乘性），但紧接着的微分里原文印作 $(\wedge_t(x)\wedge_{-t}(y))$，第一个因子应为 $\wedge_{-t}(x)$；否则下一步用 Leibniz 法则拆不出 $\Psi_t(x)+\Psi_t(y)$，两端对不上。正文照原样译，读者按 $\wedge_{-t}(x)$ 理解。]

对第二条，我们有：
\[ \wedge_{-t}(L) = 1-tL \]
因而
\[ \frac{d}{dt}\wedge_{-t}(L) = -L, \]
于是
\[ \Psi_t(L) = 1+\frac{tL}{1-tL} = 1+tL+t^2L^2\Rightarrow \Psi^k(L) = L^k. \]
% [Tutor Note: p.51 这一串是纯形式幂级数计算：$\frac{tL}{1-tL}=tL+t^2L^2+\ldots$ 是几何级数，逐项看即 $t^k$ 的系数为 $L^k$。核对点：把 $t^k$ 的系数与 $\sum t^k\Psi^k$ 对齐，得到 $\Psi^1(L)=L$、$\Psi^k(L)=L^k$（$k\ge2$），正是命题 1 第 2 条。另注：原文 "propety"（p.51 第四段）系 property 之误。]

对第四条，我们简言之：由第一条性质（可加性），若该性质对线丛成立，则由分裂原理它对任意向量丛也成立。

对第五条，设 $x = L_1+L_2+\ldots+L_n$，其中 $L_i$ 是线丛，则
\[ \Psi^p(x) = L_1^p+L_2^p+\ldots+L_n^p = (L_1+L_2+\ldots+L_n)^p modp. \]

% [原文 p.52]
对最后一条，设 $h$ 是 $K^0(S^2)$ 的生成元，$J$ 表示 $S^2 = \C P^1$ 上的 Hopf 线丛。则
\[ h = [J]-1\in \wt K^0(S^2). \]
计算
\[ \Psi^k(h) = \Psi^k[J]-\Psi^k(1) = [J^k]-1 \]
于是
\[ [J^k] = k[J]-(k-1)\in K^0(S^2). \]
利用 Chern 特征
\[ \ch : K^0(S^2) = \Z\oplus\Z\to H^{ev}(X;\Q) \]
% [Tutor Note: p.52 上式右端原文印作 $H^{ev}(X;\Q)$；这里算的是 $S^2$ 上的 Chern 特征，按上下文应为 $H^{ev}(S^2;\Q)$（同页其余各式均以 $S^2$ 为对象）。正文照原文保留 $X$。]
连同上面的关系
\[ [J^k] = k[J]-(k-1)\in K^0(S^2) \]
得到
\[ [J^k]-1 = k[J]-k. \]
然后考虑映射
\[ \wt K^0(S^2)\otimes\wt K^0(S^2)\to\wt K^0(S^4) \]
其中
\[ S^4 = S^2\wedge S^2. \]
$\wt K^0(S^{2n})$ 的生成元是 $h\otimes\ldots\otimes h$，于是
\[ \Psi^k(h\otimes\ldots\otimes h) = k^n(h\otimes\ldots\otimes h). \]
\end{proof}
% [Tutor Note: p.52 「Then the map ... where $S^4=S^2\wedge S^2$」在原文中是残句（缺主句）：它说的是外积 $\wt K^0(S^2)\otimes\wt K^0(S^2)\to\wt K^0(S^4)$，即 p.40 开头那族乘法映射在 $X=Y=S^2$ 的特例。正文照原样译，不替作者补主句。]
% [Tutor Note: p.52 末段的算式可用 Chern 特征独立核对：$\ch([J^k]-1)=k\ch([J])-k=k(1+u)-k=ku$，而 $\ch$ 是单射，故 $[J^k]-1=k([J]-1)=kh$，即 $\Psi^k(h)=kh$，与第 6 条（$n=1$）相符。$h\otimes\ldots\otimes h$ 的 $n$ 重张量再乘 $k^n$ 同理。]
